%% Chapitre 04 — Penser en tâches délégables
%% RÈGLE : uniquement des commandes sémantiques bookforge. Aucun \chapter,
%% \textbf, \begin{lstlisting}, \vspace... brut.

\chapitre{Penser en tâches délégables}

\introChapitre{
  Le vrai levier n'est pas de mieux « prompter » mais de mieux découper. Ce
  chapitre, central et durable, t'apprend à reformuler ton travail en tâches que
  l'agent peut réussir : bornées, vérifiables, avec un critère de succès
  explicite. C'est la compétence qui sépare ceux qui s'agacent de ceux qui
  gagnent du temps.
}

\section{Anatomie d'une tâche délégable}

\paragraphe{Au chapitre précédent, tu as confié une vraie tâche à l'agent et tu
l'as vu explorer, proposer, modifier. Tu as réussi un coup. La question
maintenant : comment transformer ce coup unique en réflexe, pour toutes tes
tâches ? La réponse tient dans un changement de regard. Tu arrêtes de penser
« comment lui parler » et tu commences à penser « est-ce que ce travail est
formulé pour qu'il puisse le réussir ».}

\paragraphe{C'est exactement ce qu'on appelle une \terme{tâche délégable} : un
travail reformulé pour qu'un agent puisse l'accomplir. Pas n'importe quel
travail jeté par-dessus l'épaule, mais un travail qui réunit quatre ingrédients.
Un \important{but} clair : ce qu'on veut obtenir, pas la manière. Un
\important{contexte} suffisant : ce que l'agent doit savoir et où regarder. Un
\important{périmètre} borné : ce qui est inclus et, surtout, ce qui ne l'est pas.
Et un \important{critère de succès} vérifiable : un test objectif qui dit « c'est
fini » sans dépendre de ton avis du moment.}

\paragraphe{Retiens ces quatre mots, ils structurent tout le chapitre : but,
contexte, périmètre, critère. Une demande à laquelle il manque un de ces quatre
n'est pas mauvaise, elle est simplement \important{incomplète} — et l'agent
comblera le trou en devinant, parfois bien, souvent à côté.}

\paragraphe{Prends un exemple minimal. Une mauvaise formulation : « améliore la
fonction de connexion ». Aucun des quatre ingrédients n'y est. Améliorer dans
quel sens ? Quelle fonction, dans quel fichier ? Jusqu'où va-t-on ? Et comment
saura-t-on que c'est mieux ? Compare avec une version délégable :}

\begin{extraitcode}{text}
But       : la fonction loginUser doit rejeter les mots de passe
            de moins de 8 caractères.
Contexte  : la fonction est dans src/auth/login.js ; la validation
            actuelle ne vérifie que le format de l'email.
Périmètre : ne touche qu'à loginUser ; ne change pas la signature
            ni les autres validations.
Succès    : un nouveau test rejette "abc" et accepte "motdepasse1".
\end{extraitcode}

\paragraphe{Même demande sur le fond, mais l'une se résout en une boucle, l'autre
en une devinette. Tu n'as pas besoin de réciter ces quatre lignes à chaque fois —
avec l'habitude elles tiennent dans une phrase. Ce qu'il te faut, c'est l'instinct
de vérifier qu'aucune ne manque avant d'appuyer sur Entrée.}

\begin{astuce}
Quand une de tes consignes se passe d'instructions sur la \emphase{manière}, c'est
bon signe : tu décris un résultat, pas une recette. Laisse à l'agent le « comment »
tant que tu tiens fermement le « quoi » et le « comment vérifier ».
\end{astuce}

\section{Découper une grosse demande en sous-tâches enchaînables}

\paragraphe{La plupart des échecs ne viennent pas d'une mauvaise formulation mais
d'une tâche trop grosse. « Ajoute l'authentification à mon application » n'est pas
une tâche, c'est un projet : schéma de base de données, routes, validation,
gestion de session, pages, tests. Donné d'un bloc, l'agent va soit s'éparpiller,
soit prendre dix décisions à ta place que tu découvriras trop tard dans un
\terme{diff} géant impossible à relire.}

\paragraphe{Le réflexe : découper en sous-tâches qui s'enchaînent, chacune
délégable au sens de la section précédente. Une bonne sous-tâche produit un
résultat \important{vérifiable seul}, avant de passer à la suivante. Tu ne lances
pas la deuxième tant que la première n'est pas verte.}

\begin{extraitcode}{text}
1. Crée la table `users` (email, password_hash, created_at)
   + la migration. Succès : la migration s'applique sans erreur.
2. Ajoute la route POST /signup qui valide l'email et hache
   le mot de passe. Succès : un test crée un user et renvoie 201.
3. Ajoute la route POST /login qui vérifie les identifiants.
   Succès : un test accepte le bon couple, rejette le mauvais.
4. Branche la session sur /login. Succès : un test montre
   un cookie de session posé après connexion.
\end{extraitcode}

\paragraphe{Chaque étape s'appuie sur la précédente, et chacune a son propre
critère de succès. Si l'étape 3 part de travers, tu le vois sur l'étape 3 — tu
ne remontes pas une avalanche. C'est la différence entre déboguer un changement
de cinquante lignes et déboguer un changement de huit cents.}

\paragraphe{Tu peux confier le découpage lui-même à l'agent : « voici ce que je
veux au final, propose-moi un découpage en étapes vérifiables une par une, ne
code rien encore ». Tu obtiens un plan que tu relis et corriges avant qu'une
seule ligne soit écrite. Valider le plan coûte une minute ; relire le mauvais
code en coûte trente.}

\begin{avertissement}
Une grosse demande qui « marche du premier coup » est souvent un piège : l'agent
a comblé les trous avec des hypothèses que tu n'as pas vues passer. Ça compile,
donc tu fais confiance — jusqu'à ce qu'un choix implicite te saute au visage en
production. Découper, c'est rendre ces choix visibles un par un.
\end{avertissement}

\section{Tâches à déléguer, tâches à garder pour soi}

\paragraphe{Tout n'est pas à déléguer, et croire le contraire est le meilleur
moyen de perdre confiance dans l'agent. Le bon partage ne dépend pas de la
difficulté technique mais de deux questions : \important{le critère de succès
est-il objectif}, et \important{le coût d'une erreur est-il récupérable} ?}

\paragraphe{Délègue volontiers ce qui est borné et vérifiable : écrire des tests
à partir d'une spec, renommer un symbole partout, traduire un module d'une API à
une autre, rédiger une première version de documentation, repérer où un
comportement est implémenté dans une base que tu connais mal. Ce sont des tâches
où « réussi » se constate, pas se discute.}

\paragraphe{Garde pour toi, ou ne délègue qu'après cadrage serré, ce qui repose
sur un jugement que toi seul détiens : choisir une architecture qui engage le
projet pour des mois, arbitrer un compromis produit, décider d'une stratégie de
migration risquée, trancher sur des données sensibles. Là, le critère de succès
est ton jugement — il n'est pas externalisable. Tu peux faire \emphase{explorer} les
options par l'agent, mais la décision reste à toi.}

\begin{liste}
  \element{Critère objectif + erreur récupérable → délègue sans hésiter.}
  \element{Critère objectif + erreur coûteuse → délègue, mais avec un périmètre
  étroit et une relecture serrée du \terme{diff}.}
  \element{Critère subjectif → garde la décision ; ne délègue au mieux que
  l'exploration des options.}
\end{liste}

\paragraphe{Avec le temps, ce tri devient automatique. Tu sens en lisant ta
propre demande si elle a un critère de succès que tu pourrais coder en test —
ou si elle te demande, à toi, de regarder le résultat et de juger « oui, c'est
bien comme ça ». Le premier cas se délègue ; le second se supervise.}

\section{Donner le bon contexte sans noyer l'agent}

\paragraphe{Un agent ne raisonne que sur son \terme{contexte} : ce qu'il a sous
les yeux, il l'utilise ; ce qu'il ignore, il l'invente. La tentation du débutant
est alors de tout donner — coller cinq fichiers, l'historique complet, le ticket,
la doc. Mauvais réflexe. Trop de contexte dilue les informations qui comptent et
fait dériver l'agent vers des détails sans rapport.}

\paragraphe{Le bon contexte est \important{ciblé}, pas \important{exhaustif}. La
question utile n'est pas « qu'est-ce que je pourrais lui donner ? » mais « de
quoi a-t-il strictement besoin pour réussir \emphase{cette} tâche ? ». Souvent, la
meilleure réponse n'est pas un fichier collé mais une indication d'où regarder :
l'agent ira lire lui-même ce qui lui manque, plus à jour que ton copier-coller.}

\paragraphe{Compare. Une formulation qui noie : « voici tout le dossier \texttt{src}, le README, et le ticket JIRA, débrouille-toi pour corriger le bug de date ». Une formulation qui cible : « le bug d'affichage des dates est dans \code{src/utils/format.js}, fonction \code{formatDate} ; les dates en UTC s'affichent avec un jour de retard ; regarde comment elle gère les fuseaux ». La seconde tient en une ligne et pointe l'agent droit sur la zone.}

\begin{astuce}
Préfère « regarde dans \code{tel fichier} » à « voici le contenu de \code{tel fichier} ». Donner un chemin laisse l'agent lire la version réelle et courante ;
coller du texte fige une copie qui peut déjà être périmée et qui occupe de la
place pour rien.
\end{astuce}

\paragraphe{Le contexte qui compte vraiment, c'est l'\important{implicite} : les
conventions, les contraintes, les pièges que tu as en tête et que l'agent ne peut
pas deviner. « On utilise toujours des dates en UTC en base », « ne touche jamais
au dossier \code{legacy/} », « les tests tournent avec \code{npm test} ». Cet
implicite-là vaut dix fichiers collés. C'est lui qu'il faut rendre explicite —
et c'est précisément ce que la suite du livre apprend à automatiser, pour ne pas
le répéter à chaque \terme{session}.}

\section{Spécifier le « fini » : comment l'agent sait qu'il a réussi}

\paragraphe{De tous les ingrédients, le critère de succès est celui qu'on oublie
le plus et celui qui rapporte le plus. Sans lui, l'agent s'arrête quand
\emphase{lui} estime avoir fini — c'est-à-dire au premier point qui lui semble
plausible. Avec lui, il a une cible objective : il peut vérifier son propre
travail, et toi aussi.}

\paragraphe{Le meilleur critère est \important{exécutable} : un test qui passe,
une commande qui sort sans erreur, une sortie attendue précise. « La fonction
marche » n'est pas un critère ; « \code{npm test} passe et le nouveau test
\code{rejette les emails sans @} est vert » en est un. Le premier dépend de
l'humeur, le second se constate.}

\begin{extraitcode}{text}
Flou   : « assure-toi que le parsing est robuste »
Net    : « parseConfig doit renvoyer une erreur explicite sur un
          JSON invalide ; ajoute un test qui le vérifie sur "{" »

Flou   : « rends la page plus rapide »
Net    : « le rendu initial ne doit plus déclencher qu'une seule
          requête réseau ; vérifie-le dans l'onglet réseau »

Flou   : « nettoie ce module »
Net    : « supprime les fonctions non utilisées de utils.js ;
          le lint et les tests doivent rester verts »
\end{extraitcode}

\paragraphe{Une astuce qui change tout : demande à l'agent de \important{prouver}
qu'il a réussi, pas seulement de l'affirmer. « Montre-moi le test qui le
démontre » ou « lance la commande et colle la sortie ». Tu passes de la confiance
sur parole à la preuve. Et quand un critère exécutable n'existe pas — un libellé
d'interface, un ton de message — rends-le concret autrement : un exemple d'entrée
et la sortie exacte attendue. Un exemple vaut mieux qu'un adjectif.}

\begin{astuce}
Avant de lancer une tâche, finis mentalement la phrase « je saurai que c'est
réussi quand… ». Si tu ne peux pas la terminer par quelque chose d'observable,
ta tâche n'est pas encore prête : il manque le critère, et c'est toi qui
serviras de critère à la main, à chaque essai.
\end{astuce}

\section{Anti-patterns de formulation, et comment les corriger}

\paragraphe{Les ratés se répètent. Les nommer aide à les repérer dans tes propres
demandes. En voici quatre, chacun avec sa version corrigée et le critère qui
prouve la correction.}

\paragraphe{\important{Le vœu pieux.} Tu décris un état d'esprit, pas un
résultat : « rends le code plus propre », « optimise ça », « améliore la
lisibilité ». L'agent fait quelque chose, tu ne sais pas quoi évaluer. Correction :
remplace l'adjectif par une action observable. Au lieu de « rends ce module plus
maintenable », écris « extrais la logique de validation dans une fonction
\code{validate} et ajoute un test pour chaque cas d'erreur ». Critère : la
fonction existe, les tests passent.}

\paragraphe{\important{Le périmètre fuyant.} Tu demandes une chose et tu laisses
la porte ouverte à dix autres : « corrige ce bug (et tant qu'à faire, mets à jour
ce qui traîne) ». L'agent élargit, le \terme{diff} enfle, tu ne peux plus le
relire. Correction : borne explicitement. « Corrige uniquement ce bug ; ne touche
à aucun autre fichier ; si tu vois autre chose à faire, liste-le sans le faire ».
Critère : le diff ne touche que le fichier visé.}

\paragraphe{\important{Le contexte fantôme.} Tu supposes que l'agent sait ce que
tu as en tête : « tu sais, fais comme l'autre endpoint ». Quel autre endpoint ?
L'agent invente une référence plausible. Correction : nomme la source. « Reproduis
le schéma de \code{routes/users.js}, lignes 40 à 70, pour le nouvel endpoint ».
Critère : le nouveau code suit la même structure, vérifiable côte à côte.}

\paragraphe{\important{Le test absent.} Tu décris parfaitement le travail mais tu
oublies de dire comment on vérifie qu'il est fait : l'agent s'arrête où il veut.
Correction : ajoute toujours la dernière ligne. « …et ajoute un test qui échoue
sans ta correction et passe avec ». Critère : le test existe et a bien ce
comportement — il rougit avant, verdit après.}

\begin{liste}
  \element{Vœu pieux → remplace l'adjectif par une action observable.}
  \element{Périmètre fuyant → dis ce qui est exclu, pas seulement inclus.}
  \element{Contexte fantôme → nomme la source au lieu d'y faire allusion.}
  \element{Test absent → termine par « …et prouve-le par un test ».}
\end{liste}

\paragraphe{Tu remarqueras que les quatre corrections ramènent aux quatre
ingrédients du début : but, contexte, périmètre, critère. C'est tout ce chapitre
en une boucle. Reformuler une demande agaçante en tâche délégable, ce n'est pas
un don : c'est passer ces quatre points en revue et combler celui qui manque.}

\resumeChapitre{
  \element{\important{Quatre ingrédients} : toute tâche délégable réunit un but clair, un contexte suffisant, un périmètre borné et un critère de succès vérifiable — un ingrédient manquant et l'agent devine.}
  \element{\important{Découper} : une grosse demande produit un diff géant impossible à relire ; découpe en sous-tâches dont chacune a son propre critère de succès, et valide l'une avant de lancer la suivante.}
  \element{\important{Contexte ciblé} : pointe l'agent vers les bons fichiers plutôt que de tout coller — trop d'information dilue les détails qui comptent.}
  \element{\important{Critère exécutable} : le meilleur critère est un test qui passe ou une commande sans erreur ; « ça marche » n'est pas un critère, c'est un avis.}
  \element{\important{Anti-patterns} : vœu pieux, périmètre fuyant, contexte fantôme, test absent — ces quatre défauts se corrigent chacun en rajoutant l'ingrédient qui manque.}
}

\paragraphe{Il reste un ingrédient qu'on a effleuré sans creuser : le contexte.
On a dit qu'il fallait le donner ciblé, et que l'implicite — conventions,
contraintes, pièges — comptait plus que le volume. Mais le contexte n'est pas
qu'une question de formulation ponctuelle : c'est une ressource limitée qui se
remplit, sature et se gère au fil d'une \terme{session}. C'est l'objet du
chapitre suivant.}
